%------------------------------------------------------------------------------%
\documentclass[fleqn,utf8,aspectratio=169,12pt]{beamer}

\usepackage{gaal}

\title{Espaço $\R^n$ e Suas Operações}

%------------------------------------------------------------------------------%
\begin{document}

\addtitle

%------------------------------------------------------------------------------%
\section{O espaço $\R^n$}

%------------------------------------------------------------------------------%
\begin{frame}{$\R^n$}
  Para $n\in\mathbb{N}$
  \[
    \R^n = \left\{ (x_1, x_2, \ldots, x_n) \mid x_i\in\R \right\}
  \]

  Um elemento de $\R^n$ é um vetor
  \[
    \vec{v} = (x_1, x_2, \ldots, x_n)
  \]

  Os números
  \[
    v_1, \ldots, v_n
  \]
  são as componentes de $\vec{v}$
\end{frame}

%------------------------------------------------------------------------------%
\begin{frame}{Adição em $\R^n$}
  Se
  \[
    \vec{v} = (v_1, \ldots, v_n)
  \]
  \[
    \vec{w} = (w_1, \ldots, w_n)
  \]
  então
  \[
    \vec{v} + \vec{w} = ( v_1+w_1,\, \ldots,\, v_n+w_n )
  \]
\end{frame}

%------------------------------------------------------------------------------%
\begin{frame}{Vetor Nulo}
  O vetor nulo de $\R^n$ é
  \[
    \vec{0} = (0, 0, \ldots, 0)
  \]
  \[
    \vec{v}+\vec{0} = \vec{v}
  \]
\end{frame}

%------------------------------------------------------------------------------%
\begin{frame}{Vetor Oposto}
  Se
  \[
    \vec{v} = (v_1, \ldots, v_n)
  \]
  então
  \[
    -\vec{v} = (-v_1, \ldots, -v_n)
  \]
  \[
    \vec{v}+(-\vec{v}) = \vec{0}
  \]
\end{frame}

%------------------------------------------------------------------------------%
\begin{frame}{Multiplicação por Escalar}
  Se
  \[
    \alpha\in\R
  \]
  \[
    \vec{v} = (v_1, \ldots, v_n)
  \]
  então
  \[
    \alpha\vec{v} = (\alpha v_1, \ldots, \alpha v_n)
  \]
\end{frame}

%------------------------------------------------------------------------------%
\begin{frame}{Norma em $\R^n$}
  Para
  \[
    \vec{v} = (v_1, \ldots, v_n)
  \]
  a norma é
  \[
    \norm{\vec{v}} = \sqrt{ v_1^2 + \cdots + v_n^2}
  \]
\end{frame}

%------------------------------------------------------------------------------%
\begin{frame}{Produto escalar}
  Para
  \[
    \vec{v} = (v_1, \ldots, v_n)
  \]
  \[
    \vec{w} = (w_1, \ldots, w_n)
  \]
  o produto escalar é
  \[
    \vec{v}\cdot\vec{w} = v_1w_1 + \cdots + v_nw_n
  \]
\end{frame}

%%------------------------------------------------------------------------------%
%\section{Lista Mínima}
%
%%------------------------------------------------------------------------------%
%\begin{frame}{Lista Mínima}
%  \begin{enumerate}
%    \item
%  \end{enumerate}
%
%  \vfill
%  Atenção: A prova é baseada no livro, não nas apresentações
%\end{frame}

%------------------------------------------------------------------------------%
\end{document}
%------------------------------------------------------------------------------%
